% answers_short.tex —— 2008 年·普通高等学校招生考试·新课标卷（琼、宁卷）·文科 · 简短答案速查
% 由 tools/make_short_answers.py 自 answers.tex 的 \gksolhead/\gkanswer 提取，
% 并与 2008/2008_数学真题_新课标文_答案要点.md 逐题比对后生成，勿手改（改源 MD 后重跑）。

\begin{gknotice}
\textbf{使用说明}：本节只给最终答案便于快速核对，完整推导见解析版。答案为编写者独立推导，\textbf{非官方标准答案}。
\end{gknotice}

\gkpart{选择题}
\begin{center}\begin{tabular}{ccccccccccccc}
\toprule
题号 & 1 & 2 & 3 & 4 & 5 & 6 & 7 & 8 & 9 & 10 & 11 & 12 \\
\midrule
答案 & C & B & A & B & A & A & B & C & D & B & C & D \\
\bottomrule
\end{tabular}\end{center}

\gkpart{填空题}
\begin{itemize}[leftmargin=2.6em]
  \item[\textbf{13.}] $15$。
  \item[\textbf{14.}] $\dfrac{4\pi}3$。
  \item[\textbf{15.}] $\dfrac53$。
  \item[\textbf{16.}] （开放性作答，以下两条均可）
\end{itemize}

\gkpart{解答题}
\begin{itemize}[leftmargin=2.6em]
  \item[\textbf{17.}] （1）$\cos\angle CBE=\dfrac{\sqrt6+\sqrt2}4$；（2）$AE=\sqrt6-\sqrt2$。
  \item[\textbf{18.}] （1）俯视图为 $6\times4$ 的矩形，在其左上角沿连接两邻边上距该角为 $2$ 的两点的线段切去一个等腰直角三角形；（2）$\dfrac{284}3\ \mathrm{cm^3}$；（3）证明见解析。
  \item[\textbf{19.}] （1）$7.5$；（2）$\dfrac7{15}$。
  \item[\textbf{20.}] （1）$\left[-\dfrac12,\dfrac12\right]$；（2）不能，理由见解析。
  \item[\textbf{21.}] （1）$f(x)=x-\dfrac3x$；（2）面积为定值 $6$。
\end{itemize}

\gkpart{选做题}
\begin{itemize}[leftmargin=2.6em]
  \item[\textbf{22.}] （1）（2）证明见解析。
  \item[\textbf{23.}] （1）$C_1$ 是单位圆 $x^2+y^2=1$，$C_2$ 是直线 $y=x+\sqrt2$，二者有 $1$ 个公共点；（2）$C_1':x^2+4y^2=1$（参数方程 $x=\cos\theta$，$y=\dfrac12\sin\theta$），$C_2':y=\dfrac{x+\sqrt2}2$（参数方程 $x=\dfrac{\sqrt2}2t-\sqrt2$，$y=\dfrac{\sqrt2}4t$），公共点个数相同，仍为 $1$。
\end{itemize}
